\section{Complete original curve-family smooth functoriality}

At original dimension bound one, every nonspecial maximal-order point would
produce a dimension-zero lower cosupport, which is empty. This closes the
whole first positive recursion level, rather than only a monomial or special
branch. Original histories, centres and ordered boundary entries are retained.
Sources are W\l odarczyk's dimension induction and Bierstone--Milman's smooth
functoriality \cite{wlodarczyk05-hironaka,bierstone-milman08-functoriality}.

\begin{theorem}[Original curve maximal-order canonical cosupport is special]
\label{hir:smooth-original-curve-step-one-special}
\source{wlodarczyk05-hironaka,bierstone-milman08-functoriality}
\lean{Hironaka.HironakaSmoothSubmersionOriginalCurveFamily.original_curve_step_one_special}
\leanok
\uses{hir:step-one-canonical-tail,hir:centre-function,hir:canonical-centres}
For $n>0$, original $M$ in the dimension-one domain, original maximal
order, and original Step One $R^{(n)}_0$-canonical whole history $T$, every
point of the actual final cosupport is special.
\end{theorem}
\begin{proof}
The existing Step One bridge internally supplies the true bar-marked test
history, its canonicality and the cosupport equality. If a given cosupport
point were nonspecial, the original canonical-tail theorem would supply
its genuine masked lower marked ideal, an original domain with dimension
bound $1-1=0$, its original $R^{(n)}_0$-canonical restricted history, and
membership of that same point in the lower transformed cosupport. The
zero-domain theorem makes that cosupport empty, a contradiction. This
handles both positive-member and zero-contact branches without comparing
independently chosen contact functions or assuming a lower certificate.
\end{proof}

\begin{theorem}[Complete actual original curve-family smooth values]
\label{hir:smooth-original-curve-full-family-values}
\source{wlodarczyk05-hironaka,bierstone-milman08-functoriality}
\lean{Hironaka.HironakaSmoothSubmersionOriginalCurveFamily.IsActualTowerPullback.actual_curve_family_value}
\leanok
\uses{hir:smooth-original-curve-step-one-special,hir:smooth-original-family-monomial-leaf,hir:smooth-actual-value-recursive-dimension-maps,hir:submersion-relative-control-joint-tower,hir:general-canonical-tail}
Let $n,m>0$, $t=m-n$, and let an actual smooth Cartesian pair $(Q,T)$
pull back the original dimension-one domain $M$ and its original whole
$R^{(n)}_1$-canonical history. At every $z$ whose image lies in the actual
target cosupport, the complete actual original family values satisfy
\[R^{(m)}_{t+1}(f^*M,Q,z)=D_{t,1}(R^{(n)}_1(M,T,gz)).\]
Only original target canonicality is assumed; source canonicality and a
source value formula are not premises.
\end{theorem}
\begin{proof}
Target canonicality gives its true original test history. If its actual
nonmonomial order at the image is zero, consume the established complete
smooth monomial-leaf formula at depth one. Otherwise write the actual
order as a positive integer $o$. The original general canonical-tail
theorem, with the genuine Step One maximal-order value over $R^{(n)}_0$,
internally supplies the selected order-mask membership, whole companion
domain of dimension one, maximal order, its whole Step One canonical
history and its actual cosupport point. The first theorem makes this
companion point special. The existing actual historical-tail constructor
gives the literal whole source companion record, both restricted smooth
maps, the actual whole Cartesian restricted histories and their point
identity. Original specialness transports by the actual geometric
dictionary, and its full Step One value is the literal pair
$(\mathrm{level},\top)$ on both sides, with the same actual level. Thus
the established strict Step One dimension map compares both entire
companion values. The actual order and run dictionaries give the same
outer $o$ and exact split index; expanding the original general-value
definition then yields precisely $D_{t,1}$ of the whole target value.
No positive-inner or independent-contact case is left at this true depth.
\end{proof}

\begin{theorem}[Complete original curve-family canonical history pullback]
\label{hir:smooth-original-curve-full-family-canonical}
\source{wlodarczyk05-hironaka,bierstone-milman08-functoriality}
\lean{Hironaka.HironakaSmoothSubmersionOriginalCurveFamily.IsActualTowerPullback.original_curve_family_canonical}
\leanok
\uses{hir:smooth-original-curve-full-family-values,hir:smooth-actual-value-recursive-dimension-maps,hir:submersion-relative-control-joint-tower}
Under precisely the preceding original inputs, the actual whole source
history $Q$ is canonical for the actual original $R^{(m)}_{t+1}$, retaining
every empty step and every literal preimage centre.
\end{theorem}
\begin{proof}
Induct on the actual Cartesian history witness. The trivial pair is
canonical. The original last centre is either empty or the full maximum
locus of the preceding target history. The induction hypothesis supplies
source canonicality of the preceding history. An empty target centre or
an empty source preimage yields an allowed empty source step. Otherwise
choose a source point mapping to the actual full target maximum. The
actual cosupport dictionary and the preceding complete value theorem,
applied at every source cosupport point, show that the source maximum
locus is exactly the preimage of that full target maximum. Indeed, a
source maximum dominates the chosen visible target maximum by strict
order reflection, which makes its image a full target maximum; the
reverse implication follows from domination of all target points.
The strict unchanged $D_{t,1}$ supplies both order implications. The
literal source centre is therefore its own full maximum locus. No
surjectivity, source compactness or global-maximum existence is assumed,
and the required actual value formulas are proved internally at every
inductive stage.
\end{proof}

\begin{theorem}[Original compact curve input constructs both terminal histories]
\label{hir:smooth-original-curve-full-family-resolution}
\source{wlodarczyk05-hironaka,bierstone-milman08-functoriality}
\lean{Hironaka.HironakaSmoothSubmersionOriginalCurveFamily.exists_original_smooth_curve_family_resolution}
\leanok
\uses{hir:smooth-original-curve-full-family-canonical,hir:invariant-driver-seq,hir:submersion-relative-invariant-domain,hir:submersion-relative-control-joint-tower}
For compact original target $X$, original smooth $f:Y\to X$, $n,m>0$ and
original dimension-one domain $M$, internally construct a target whole
canonical terminal history $T$ for $R^{(n)}_1$ and its actual whole source
$Q,g$. The source has the correct domain bound $t+1$, a whole canonical
terminal history for the actual original $R^{(m)}_{t+1}$, unchanged history
length, and the literal whole marked and cosupport pullback dictionaries.
\end{theorem}
\begin{proof}
The existing per-sequence invariant driver applied to the original
$R^{(n)}_1$ and original $M$ supplies target canonicality and terminal
geometric resolution, using compactness only of $X$. The actual smooth
tower constructor supplies $Q,g$ and its entire Cartesian witness.
Apply the preceding theorem for source canonicity. The original-domain
pullback theorem supplies the true source bound $1+t=t+1$ and all other
fields internally. The actual marked dictionary supplies source testness
and whole transformed-record equality. The cosupport dictionary makes
the final source cosupport the preimage of the target's empty final
cosupport, giving source terminal resolution. The actual witness supplies
unchanged length. This closes the whole original bound-one smooth family
case; arbitrary higher original dimension remains an active obligation.
\end{proof}
